\documentclass[11pt]{article}

\usepackage[margin=1in]{geometry}
\usepackage{amsmath,amssymb,amsthm,mathtools}
\usepackage{booktabs}
\usepackage{microtype}
\usepackage{hyperref}

\hypersetup{colorlinks=true,linkcolor=blue,citecolor=blue,urlcolor=blue}
\newtheorem{theorem}{Theorem}
\newtheorem{lemma}{Lemma}
\newtheorem{corollary}{Corollary}
\theoremstyle{definition}
\newtheorem{definition}{Definition}
\newtheorem{remark}{Remark}
\newcommand{\HH}{\mathcal H}
\newcommand{\dd}{\,\mathrm d}
\newcommand{\norm}[1]{\left\lVert #1\right\rVert}
\newcommand{\ip}[2]{\left\langle #1,#2\right\rangle}
\newcommand{\Vhat}{\widehat V}
\newcommand{\ii}{\mathrm i}

\title{A Null-Compatible Four-Packet Completion Obstruction\\
for the Literal V59 Interface}
\author{Liang Wang\\
School of Mathematics and Statistics, Huazhong University of Science and Technology\\
Wuhan, China}
\date{August 26, 2026}

\begin{document}
\maketitle

\begin{abstract}
The preceding same-clock analysis suppresses one source-frozen rank-one channel of
the signed V59 coupling, but leaves an orthogonal residual.  We give an exact finite
audit of whether four packet marginals and the existing Haar information can control
that residual.  First, a polygon-completion theorem gives the sharp range of the
full residual for fixed packet norms.  Second, a four-point discrete Fourier
transform (DFT) identifies the missing datum: packet diagonal energy fixes the sum
of all mode energies, whereas the full reassembly is precisely mode zero.  In the
equal-norm case, two packet families have identical norms, vanish on every Haar
contrast and on the TPC-259 null channel, yet have full residual energies $16$ and
$0$.  The result is a structural obstruction, not a prime-shell counterexample.
It converts the next arithmetic task into a precise one: prove a common-clock
mode-zero or cross-Gram estimate while retaining the hard window, deleted diagonal,
unit masks, and the physical residual.
\end{abstract}

\section{Scope and motivation}

The V59 Route-B object is a signed coupling on a fixed real clock.  Recent work
constructed a four-block Haar frame and then chose a source-frozen transverse vector
that cancels the known diagonal asymptotic.  Coupling that vector to the literal
hybrid sequence gives an exact split
\begin{equation}
 \ip{w}{A_x\beta}
 =\overline{\ip{z_{\rm null}}{w}}\,\ip{z_{\rm null}}{A_x\beta}
   +\ip{w_\perp}{A_x\beta},
 \qquad
 w_\perp=w-\ip{z_{\rm null}}{w}z_{\rm null}.
 \label{eq:tpc259-split}
\end{equation}
The first term is source-backed and arbitrarily small in fixed logarithmic powers;
the second term is not estimated.  This paper asks a deliberately narrower
question: can the residual be recovered from packet norms and the finite Haar
projections already present in the interface?

The answer is no at the level of finite Hilbert-space algebra.  The obstruction is
useful because it does not merely repeat a generic Gram example.  It keeps the
specific source-frozen null direction, embeds it in the four-block Haar complement,
and gives a sharp continuum of compatible residual sizes.  Thus it tells us exactly
which information a future literal theorem must add.

Our contributions are:
\begin{enumerate}
 \item an exact weighted four-block Haar complement containing the TPC-258 null
       direction;
 \item a sharp polygon interval for the residual under fixed packet marginals and
       a zero null channel;
 \item a four-packet DFT ledger showing that the unresolved quantity is mode zero,
       not the total diagonal energy; and
 \item an exact plus/alternating witness, checked over varied block sizes, that
       refutes uniform promotion from these marginals to full reassembly.
\end{enumerate}

All statements below are finite and structural.  In particular, no finite record
is used as evidence for a growing prime-distribution estimate.

\section{The four-block interface}

Let $B_0,B_1,B_2,B_3$ be consecutive nonempty blocks with lengths
$s_0,s_1,s_2,s_3$ and $N=s_0+s_1+s_2+s_3$.  On blockwise-constant vectors use
the weighted inner product
\[
 \ip{u}{v}_s=\sum_{j=0}^3s_j\,\overline{u_j}v_j.
\]
Write $\ell=s_0+s_1$ and $r=s_2+s_3$.  The three contrasts used in the recent
V59 construction are
\begin{align*}
 z_0&=\rho_0\left(\frac1\ell,\frac1\ell,-\frac1r,-\frac1r\right),
 &\rho_0^2&=\frac{\ell r}{N},\\
 z_1&=\rho_1\left(\frac1{s_0},-\frac1{s_1},0,0\right),
 &\rho_1^2&=\frac{s_0s_1}{\ell},\\
 z_2&=\rho_2\left(0,0,\frac1{s_2},-\frac1{s_3}\right),
 &\rho_2^2&=\frac{s_2s_3}{r}.
\end{align*}
Direct counting gives
\begin{equation}
 \ip{z_a}{z_b}_s=\delta_{ab},\qquad
 \ip{\boldsymbol 1}{z_a}_s=0 \quad (0\leq a,b\leq2).
 \label{eq:haar}
\end{equation}
The normalized scaling direction is $e_{\rm sc}=\boldsymbol 1/\sqrt N$.
Consequently, $z_0,z_1,z_2,e_{\rm sc}$ form an orthonormal basis of the
four-dimensional block space.  In particular, the source-frozen vector
\begin{equation}
 z_{\rm null}=\frac{L_2z_1-L_1z_2}{\sqrt{L_1^2+L_2^2}},\qquad
 L_1=\log\frac{3456}{3125},\quad
 L_2=\log\frac{884736}{823543},
 \label{eq:null}
\end{equation}
is orthogonal to $e_{\rm sc}$.  Notice that this conclusion uses only the exact
zero weighted sums in \eqref{eq:haar}; it does not use numerical values for the
logarithms.

For comparison, the prior four-packet polarization result established that a
general packet Gram matrix is positive semidefinite and that diagonal/trace data
do not determine signed cross terms.  Here we specialize the obstruction to the
new same-clock residual: the packets will lie in the scaling complement of the
Haar contrasts, so all three contrast measurements are fixed at zero.

\section{A sharp null-compatible completion theorem}

Let $z,w$ be orthonormal vectors in a complex Hilbert space $\HH$.  We regard $z$
as the source-frozen null direction and $w$ as a possible residual direction.  A
four-packet family is $V_0,V_1,V_2,V_3\in\HH$ and its full output is
$G=\sum_jV_j$.

\begin{theorem}[polygon completion]
Fix nonnegative packet lengths $d_0,d_1,d_2,d_3$.  Among the families
\[
 V_j=d_j e^{\ii\theta_j}w,
 \qquad \theta_j\in\mathbb R,
 \]
the null condition $\ip{z}{G}=0$ holds identically.  If
\[
 D=\sum_{j=0}^3d_j,\qquad d_{\max}=\max_jd_j,
\]
then the possible residual moduli
\[
 R=\ip{w}{G}
\]
are exactly
\begin{equation}
 \max\{2d_{\max}-D,0\}\leq |R|\leq D.
 \label{eq:polygon}
\end{equation}
Every modulus in this interval and every argument is attained.
\end{theorem}

\begin{proof}
The scalar $R=\sum_jd_je^{\ii\theta_j}$ is the sum of four planar vectors with
prescribed lengths.  The triangle inequality gives the upper endpoint.  If one
side is longer than the other three combined, reversing the shorter sides gives
the lower endpoint $d_{\max}-(D-d_{\max})$.  Otherwise the four sides close to a
polygon, so the lower endpoint is zero.  Rotating one side continuously fills the
interval between the two endpoints; a common rotation supplies any argument.
Since $z\perp w$, every member of the family has $\ip{z}{G}=0$.  The same two
triangle inequalities apply to every choice of phases, proving sharpness.
\end{proof}

For equal packet lengths $d_j=1$, the interval is $[0,4]$, and the squared
residual can range over $[0,16]$.  This is a completion statement, not merely two
isolated examples.

\begin{corollary}[Haar-complement realization]
In the four-block space above, take $z=z_{\rm null}$ and $w=e_{\rm sc}$.  The
families in the theorem have zero projection on $z_0,z_1,z_2$, zero TPC-259
rank-one coefficient $\ip{z}{w}$, and residual range \eqref{eq:polygon}.
\end{corollary}

\begin{proof}
Equation \eqref{eq:haar} makes $e_{\rm sc}$ orthogonal to each contrast, and
\eqref{eq:null} makes it orthogonal to $z_{\rm null}$.  Substitute these vectors
into the theorem.
\end{proof}

The corollary is the key distinction from a purely abstract diagonal counterexample:
the same finite Haar geometry used to select the null direction supplies a concrete
orthogonal completion direction.  It also makes clear why suppressing the first
term in \eqref{eq:tpc259-split} says nothing by itself about the second.

\section{The missing DFT mode}

The polygon theorem describes the size ambiguity.  The following identity locates
its information-theoretic source.  Define the four-point packet DFT by
\begin{equation}
 \Vhat_k=\frac12\sum_{j=0}^3\ii^{-jk}V_j,
 \qquad 0\leq k\leq3.
 \label{eq:dft}
\end{equation}

\begin{theorem}[mode ledger]
For arbitrary packets in a complex Hilbert space,
\begin{align}
 \sum_{k=0}^3\norm{\Vhat_k}^2&=\sum_{j=0}^3\norm{V_j}^2, \label{eq:parseval}\\
 G=\sum_{j=0}^3V_j&=2\Vhat_0, \label{eq:inverse}\\
 \norm{G}^2&=4\norm{\Vhat_0}^2. \label{eq:modezero}
\end{align}
Thus packet diagonal energy determines the sum of mode energies, while full
reassembly asks for the individual mode-zero energy.
\end{theorem}

\begin{proof}
The matrix $(2^{-1}\ii^{-jk})_{k,j}$ is unitary, which gives
\eqref{eq:parseval}.  Fourier inversion at index zero gives \eqref{eq:inverse},
and taking norms gives \eqref{eq:modezero}.
\end{proof}

The distinction is exact in the two families
\begin{align*}
 V_j^+&=w, & (j=0,1,2,3),\\
 V_j^-&=(-1)^jw, & (j=0,1,2,3).
\end{align*}
Both have packet diagonal $(1,1,1,1)$, total packet energy $4$, zero projection
on every Haar contrast, and zero null channel.  Yet
\[
 \bigl(\norm{\Vhat_k^+}^2\bigr)_{k=0}^3=(4,0,0,0),
 \qquad
 \bigl(\norm{\Vhat_k^-}^2\bigr)_{k=0}^3=(0,0,4,0),
\]
so their full residual energies are $16$ and $0$.  The alternating family has
simply moved all energy from mode zero to mode two.  A fourth-root rotating family
similarly moves it to mode one.

The full energy can also be written directly as
\begin{equation}
 \norm{G}^2=\sum_j\norm{V_j}^2
  +2\operatorname{Re}\sum_{0\leq j<\ell\leq3}\ip{V_j}{V_\ell}.
 \label{eq:crossgram}
\end{equation}
Therefore a literal proof must estimate a phase-sensitive cross-Gram sum (or an
equivalent mode-zero quantity).  A bound on the diagonal sum alone cannot do so.

\section{Finite audit and route consequence}

The release certificate reconstructs 128 varied positive four-block families.
For each family it checks all three exact Haar norms, all three pairwise
orthogonalities, all three scaling complements, and the null/scaling complement.
It then checks the plus, alternating, and fourth-root phase patterns using exact
Gaussian rational arithmetic.  The resulting counts are shown in Table~\ref{tab:checks}.

\begin{table}[t]
\centering
\caption{Finite exact audit counts.  These are reproducibility checks, not
asymptotic evidence.}
\label{tab:checks}
\begin{tabular}{lr}
\toprule
Quantity & Count\\
\midrule
Four-block families & 128\\
Haar norm identities & 384\\
Haar orthogonality identities & 384\\
Scaling-complement identities & 384\\
Null/scaling complements & 128\\
Phase families & 3\\
Exact endpoint residual energies & $0,16$\\
\bottomrule
\end{tabular}
\end{table}

The producer, independent checker, and stress checker all pass in both ordinary
and optimized Python modes.  The PDF and certificate are frozen by the accompanying
bridge checker.  These checks establish the algebra and guard against accidental
promotion of a finite witness to a prime theorem; they do not estimate the literal
V59 mode-zero output.

The resulting claim firewall is:
\begin{center}
\begin{tabular}{ll}
\toprule
Item & Status\\
\midrule
Haar complement and polygon completion & PROVED\_EXACT\_FINITE\\
Four-packet DFT ledger & PROVED\_EXACT\\
Null-compatible residual non-identifiability & REFUTED\_SCOPED\\
Literal prime-shell counterexample & NONE\\
Arithmetic $L^2$ and full Gate B & NONE / OPEN\\
Strict global $1/400$ and fixed atom & UNPAID / ZERO\\
Twin-prime conclusion & NONE\\
\bottomrule
\end{tabular}
\end{center}

This is a scoped Route-B advance: it removes a tempting but invalid promotion rule
and leaves a sharply stated positive target.  The next theorem should control
\[
 \left|\ip{w_\perp}{\sum_{j=0}^3V_j}\right|
 \quad\text{or equivalently the mode-zero/cross-Gram quantity in
 \eqref{eq:modezero}--\eqref{eq:crossgram},}
\]
for the common literal V59 packet output.  It must retain the hard window, deleted
diagonal, unit masks, and both boundary lanes.  The finite completion result gives
no license to drop any of them.

\section{Conclusion}

TPC-259's source-backed null-channel suppression is a genuine projected result, but
the orthogonal residual is an independent degree of freedom.  TPC-260 makes this
statement exact: fixed packet marginals and all existing four-block contrast data
allow a sharp polygon range, including residual energies zero and sixteen.  A DFT
rephrasing shows that the missing arithmetic information is mode selection, or the
equivalent signed cross-Gram correlation.

The result is intentionally modest.  It neither proves nor disproves a cancellation
in the growing prime shell, and it pays no arithmetic $L^2$ or endpoint credit.  Its
value is diagnostic: the next useful paper must attack a common-clock mode-zero
estimate directly.  The local Route-A/Route-B evaluator files named in the Session
plan are absent from this checkout; the proof package, theorem ledger, certificate,
bridge checker, and repository policy are the fail-closed evaluation record.

\bibliographystyle{plain}
\begin{thebibliography}{9}
\bibitem{tpc222}
L.~Wang.
\newblock Four-packet polarization and the PSD cross-term obstruction.
\newblock Session report TPC-222, 2026.

\bibitem{tpc257}
L.~Wang.
\newblock Four-block Haar lifts and a transverse norm floor for the literal V59 adjoint.
\newblock Session report TPC-257, 2026.

\bibitem{tpc258}
L.~Wang.
\newblock A source-frozen transverse null direction for the literal V59 adjoint.
\newblock Session report TPC-258, 2026.

\bibitem{tpc259}
L.~Wang.
\newblock A same-clock null-channel decomposition for the literal V59 signed coupling.
\newblock Session report TPC-259, 2026.
\end{thebibliography}

\end{document}
